\subsection{Самостоятельное исследование гиперпараметров}
Перед запуском сформулируйте ожидаемое влияние каждого изменения.
Сравнивайте варианты с исходной конфигурацией, меняя один фактор за раз.
Для каждого запуска создавайте новую сеть и новый оптимизатор.
Продолжение обучения старой сети не является независимым опытом.

\begin{keybox}{Фиксированные условия}
Разделение 80/20; seed 40; одинаковая подготовка признаков и цели;
ReLU в скрытых слоях; линейный выход; скорость обучения \(10^{-3}\).
Основные опыты проводятся без раннего останова.
\end{keybox}

Таблица задаёт компактный план из десяти конфигураций. Это предложенный
учебный протокол, а не результаты подбора оптимальной модели.
Преподаватель может скорректировать бюджет до начала занятия.

\begin{table}[H]
\centering\small
\caption{План сравнения. В столбце «слои» указаны ширины скрытых слоёв.}
\begin{tabular}{@{}llllll@{}}
\toprule
Опыт & Слои & Эпохи & Пакет & Оптимизатор & Потеря \\
\midrule
B0 & 64 & 100 & 32 & Adam & MSE \\
W1 & 32 & 100 & 32 & Adam & MSE \\
W2 & 128 & 100 & 32 & Adam & MSE \\
D1 & 64, 32 & 100 & 32 & Adam & MSE \\
E1 & 64 & 30 & 32 & Adam & MSE \\
E2 & 64 & 200 & 32 & Adam & MSE \\
P1 & 64 & 100 & 16 & Adam & MSE \\
P2 & 64 & 100 & 64 & Adam & MSE \\
O1 & 64 & 100 & 32 & SGD & MSE \\
L1 & 64 & 100 & 32 & Adam & MAE \\
\bottomrule
\end{tabular}
\end{table}

\begin{enumerate}
\item \textbf{Ширина слоя.} Сравните B0, W1 и W2. Как изменились число
параметров, ошибка и разрыв между обучением и валидацией?
\item \textbf{Глубина.} Сравните B0 и D1. Объясните, какие изменения
вносит второй скрытый слой и оправдалось ли усложнение.
\item \textbf{Число эпох.} Сравните E1, B0 и E2. Покажите кривые
и определите, продолжалось ли улучшение к концу обучения.
\item \textbf{Размер мини-выборки.} Сравните P1, B0 и P2.
Рассчитайте число обновлений весов для каждого варианта и объясните,
почему одинаковое число эпох не означает одинаковый вычислительный опыт.
\item \textbf{Оптимизатор.} Сравните B0 и O1 при фиксированном шаге.
Вывод ограничьте этими настройками: он не доказывает общего превосходства
Adam или SGD при индивидуально подобранных скоростях обучения.
\item \textbf{Функция потерь.} Сравните B0 и L1 по одинаковым MAE
и RMSE в единицах цены. Разберите изменение крупных ошибок.
\end{enumerate}

\subsubsection*{Журнал результатов}
Для каждого опыта сохраните идентификатор, архитектуру, число параметров,
заданные и выполненные эпохи, размер пакета, оптимизатор, скорость
обучения, функцию потерь, seed, время обучения, MAE, RMSE и при
необходимости \(R^2\). Пример строки ниже содержит названия полей,
а не вымышленные результаты.

\begin{lstlisting}
row = {
    "experiment": "B0",
    **result["config"],
    **result["metrics"],
}
pd.DataFrame([row]).to_csv("experiment_B0.csv", index=False)
\end{lstlisting}

\subsubsection*{Как выбирать конфигурацию}
Сначала сравните MAE валидации, затем рассмотрите RMSE, сложность и
диагностические графики. Если разница мала, не объявляйте её устойчивым
преимуществом по одному запуску. Зафиксируйте, что обнаружено на данном
разбиении, и какая дополнительная проверка нужна.

Гиперпараметры не подбираются по \file{test.csv}: в нём нет правильных
ответов. В этой работе нет заданного универсального порога MAE или
требования обязательно превзойти все другие алгоритмы. Оценивается
корректность эксперимента и обоснованность выводов.

\subsubsection*{Дополнительные исследования}
По согласованию с преподавателем выполните одно расширение:
\begin{itemize}
\item включите ранний останов по MAE валидации и сравните с B0;
укажите фактическую длину обучения и правило возврата весов;
\item исследуйте влияние отдельной скорости обучения у SGD;
\item сравните фиксированное удаление пяти признаков с сохранением
их содержательно интерпретированных категорий;
\item повторите две выбранные конфигурации с несколькими заранее
заданными seed и покажите разброс на том же фиксированном разбиении.
\end{itemize}
Не смешивайте дополнительные опыты с исходной таблицей без указания,
какие условия изменились.

\clearpage
\subsection{Типичные ошибки и способы проверки}
\small
\begin{longtable}{@{}>{\raggedright\arraybackslash}p{0.27\textwidth}>{\raggedright\arraybackslash}p{0.30\textwidth}>{\raggedright\arraybackslash}p{0.35\textwidth}@{}}
\toprule
Наблюдение & Возможная причина & Проверка \\
\midrule\endhead
Разное число входных признаков & Кодировщик обучен отдельно на каждой
таблице & Один \file{fit} на обучении; одинаковые имена и порядок столбцов. \\
Подозрительно малая ошибка & Целевая цена попала в признаки или
статистики получены до разделения & Проверить список входов и место
вызова \file{fit}. \\
Ошибка обучения равна NaN & Остались NaN, бесконечности или обучение
расходится & Проверить массивы и цель; затем исследовать масштаб и шаг. \\
Прогнозы около нуля & Забыто обратное преобразование стандартизованной
цели & Применить тот же \file{target\_scaler}, который обучен на \(y_{train}\). \\
Плохое качество у дорогих домов & Большие ошибки на части диапазона,
редкие сочетания признаков & Изучить остатки и крупнейшие промахи;
не удалять их из валидации ради метрики. \\
Отрицательные исходные цены & Ошибочная шкала или неподходящее
поведение модели & Проверить обратное преобразование; обосновать
изменение модели или постобработку на валидации. \\
Результат меняется при повторе & Другое разбиение, случайная
инициализация, продолжение обучения & Сопоставить seed, версии,
идентификаторы строк и создание новой модели. \\
Ошибка размера тензоров в PyTorch & Цель имеет форму \((b,)\),
выход --- \((b,1)\) & Явно привести обе формы к \((b,1)\). \\
\bottomrule
\end{longtable}
\normalsize
